\documentclass[../../main.tex]{subfiles}
\begin{document}

\paragraph{【新課程】}
$
\begin{pmatrix} 0 & 0 \\ 0 & 0
\end{pmatrix}=O$，$
\begin{pmatrix} 1 & 0 \\ 0 & 1
\end{pmatrix}=I$とかく．
実数$a$，$b$に対し$A=
\begin{pmatrix} 0 & a \\ 1 & b
\end{pmatrix}$とおく．
いま$xI+yA$（$x$，$y$は実数）の形に表わされる行列全体からなる集合を$R$とし，
$R$から$O$を除いた集合を$G$とする．
\begin{enumerate}[(i)]
  \item $R$に属する任意の$2$つの行列の積は$R$に属することを示せ．
  \item $G$に属する任意の行列が逆行列をもつとき，
    点$(a,b)$はどのような範囲にあるか．
    これを図示せよ．
\end{enumerate}

\paragraph{【旧課程】}
$a$，$b$，$c$，$d$を実数として$f(x)=x^4+ax^3+bx^2+cx+d$とおく．
\begin{enumerate}[(i)]
  \item 方程式$f(x)=0$が$4$個の相異なる実根をもつとき，
    実数$k$に対して，方程式$f(x)+kf'(x)=0$の実根の個数を求めよ．
  \item $2$つの方程式$f(x)=0$，$f''(x)=0$が$2$個の相異なる実根を共有するとき，
    曲線$y=f(x)$は$y$軸に平行なある直線に関して対称であることを示せ．
\end{enumerate}

\end{document}
